\chapter{Impl\'ementation du backend}

Ce chapitre d\'eveloppe en d\'etail l'impl\'ementation du serveur backend du syst\`eme Andon, en pr\'esentant la structure du code, les middleware, les routes API et les services m\'etier.

\section{Point d'entr\'ee : server.js}

\subsection{Initialisation du serveur}

Le fichier \texttt{server.js} constitue le point d'entr\'ee principal. Il configure et initialise l'ensemble des composants :

\begin{lstlisting}[language=JavaScript,caption={Initialisation du serveur Express},label={lst:init_express}]
const express = require('express');
const cors = require('cors');
const http = require('http');
const { Server } = require('socket.io');
const mqtt = require('mqtt');
const { Pool } = require('pg');

const app = express();
const server = http.createServer(app);

// Middleware
app.use(cors());
app.use(express.json());

// Socket.IO
const io = new Server(server, {
  cors: { origin: '*' }
});

// PostgreSQL
const pool = new Pool({
  host: process.env.DB_HOST,
  port: process.env.DB_PORT,
  database: process.env.DB_NAME,
  user: process.env.DB_USER,
  password: process.env.DB_PASS,
  max: 20
});
\end{lstlisting}

\subsection{Connexion MQTT et gestion des messages}

\begin{lstlisting}[language=JavaScript,caption={Gestion des messages MQTT},label={lst:mqtt_handler}]
const mqttClient = mqtt.connect(
  process.env.MQTT_URL, {
    clientId: 'sews-server',
    clean: true,
    reconnectPeriod: 5000
  }
);

mqttClient.on('connect', () => {
  console.log('MQTT connecte');
  mqttClient.subscribe('sews/andon/+/alert');
  mqttClient.subscribe('sews/capteurs/+/status');
});

mqttClient.on('message', async (topic, msg) => {
  const parts = topic.split('/');
  const type = parts[1];
  const ligneId = parts[2];
  const data = JSON.parse(msg.toString());

  if (type === 'andon') {
    await handleAndonAlert(ligneId, data);
  } else if (type === 'capteurs') {
    await handleCapteurStatus(ligneId, data);
  }
});
\end{lstlisting}

\subsection{Traitement des alertes Andon}

\begin{lstlisting}[language=JavaScript,caption={Traitement d'une alerte Andon},
label={lst:handle_alert}]
async function handleAndonAlert(ligneId, data) {
  const query = `
    INSERT INTO evenements 
      (ligne_id, type, description, date_debut)
    VALUES ($1, $2, $3, NOW())
    RETURNING *
  `;
  const values = [
    ligneId,
    data.type || 'panne',
    data.description || ''
  ];

  const result = await pool.query(query, values);
  const evenement = result.rows[0];

  // Diffusion Socket.IO
  io.to('dashboard').emit('andon-event', {
    evenement,
    ligneId: parseInt(ligneId)
  });

  // Commande LED rouge
  const cmdTopic = `sews/command/${ligneId}/led`;
  mqttClient.publish(cmdTopic, 
    JSON.stringify({ rouge: true, verte: false })
  );
}
\end{lstlisting}

\section{Middleware d'authentification}

\subsection{G\'en\'eration et validation des tokens JWT}

\begin{lstlisting}[language=JavaScript,caption={Middleware JWT},label={lst:middleware_jwt}]
const jwt = require('jsonwebtoken');
const SECRET = process.env.JWT_SECRET;

function genererToken(user) {
  return jwt.sign(
    { id: user.id, role: user.role },
    SECRET,
    { expiresIn: '24h' }
  );
}

function authMiddleware(req, res, next) {
  const token = req.headers.authorization
    ?.split(' ')[1];

  if (!token) {
    return res.status(401).json({
      error: 'Token requis'
    });
  }

  try {
    const decoded = jwt.verify(token, SECRET);
    req.user = decoded;
    next();
  } catch (err) {
    return res.status(401).json({
      error: 'Token invalide'
    });
  }
}
\end{lstlisting}

\section{Routes API}

\subsection{Route d'authentification}

\begin{lstlisting}[language=JavaScript,caption={Route de connexion},label={lst:route_login}]
const bcrypt = require('bcrypt');

app.post('/api/auth/login', async (req, res) => {
  const { email, password } = req.body;

  const result = await pool.query(
    'SELECT * FROM utilisateurs WHERE email=$1',
    [email]
  );

  if (result.rows.length === 0) {
    return res.status(401).json({
      error: 'Identifiants incorrects'
    });
  }

  const user = result.rows[0];
  const valid = await bcrypt.compare(
    password, user.mot_de_passe
  );

  if (!valid) {
    return res.status(401).json({
      error: 'Identifiants incorrects'
    });
  }

  const token = genererToken(user);
  res.json({ token, user: {
    id: user.id,
    nom: user.nom,
    prenom: user.prenom,
    role: user.role
  }});
});
\end{lstlisting}

\subsection{Routes CRUD des \'ev\'enements}

\begin{lstlisting}[language=JavaScript,caption={Routes des \'ev\'enements},label={lst:routes_evenements}]
// GET /api/evenements
app.get('/api/evenements', 
  authMiddleware, async (req, res) => {
  const { page = 1, limite = 20 } = req.query;
  const offset = (page - 1) * limite;

  const result = await pool.query(`
    SELECT e.*, l.nom as ligne_nom
    FROM evenements e
    JOIN lignes l ON e.ligne_id = l.id
    ORDER BY e.date_debut DESC
    LIMIT $1 OFFSET $2
  `, [limite, offset]);

  res.json(result.rows);
});

// PUT /api/evenements/:id
app.put('/api/evenements/:id', 
  authMiddleware, async (req, res) => {
  const { statut, compte_rendu } = req.body;

  const result = await pool.query(`
    UPDATE evenements 
    SET statut = $1, 
        compte_rendu = $2,
        date_fin = CASE WHEN $1='resolu' 
                   THEN NOW() ELSE date_fin END
    WHERE id = $3
    RETURNING *
  `, [statut, compte_rendu, req.params.id]);

  // Diffusion Socket.IO
  io.to('dashboard').emit('resolution', {
    evenement: result.rows[0]
  });

  // Commande LED verte
  const ligneId = result.rows[0].ligne_id;
  mqttClient.publish(
    `sews/command/${ligneId}/led`,
    JSON.stringify({ rouge: false, verte: true })
  );

  res.json(result.rows[0]);
});
\end{lstlisting}

\subsection{Route de rapports}

\begin{lstlisting}[language=JavaScript,caption={Route de rapports},
label={lst:route_rapports}]
app.get('/api/rapports', 
  authMiddleware, async (req, res) => {
  const { dateDebut, dateFin } = req.query;

  const stats = await pool.query(`
    SELECT 
      l.nom as ligne,
      COUNT(*) as nb_evenements,
      AVG(EXTRACT(EPOCH FROM 
        (COALESCE(date_fin, NOW()) 
         - date_debut))) as duree_moyenne,
      COUNT(CASE WHEN statut='resolu' 
            THEN 1 END) as nb_resolus
    FROM evenements e
    JOIN lignes l ON e.ligne_id = l.id
    WHERE e.date_debut BETWEEN $1 AND $2
    GROUP BY l.nom
    ORDER BY nb_evenements DESC
  `, [dateDebut, dateFin]);

  res.json(stats.rows);
});
\end{lstlisting}

\section{Gestion des notifications}

\subsection{Notifications par e-mail}

Le serveur int\`egre un module d'envoi d'e-mails via Nodemailer :

\begin{lstlisting}[language=JavaScript,caption={Envoi de notification e-mail},
label={lst:send_email}]
const nodemailer = require('nodemailer');

const transporter = nodemailer.createTransport({
  host: 'smtp.sews.ma',
  port: 587,
  secure: false,
  auth: {
    user: process.env.SMTP_USER,
    pass: process.env.SMTP_PASS
  }
});

async function envoyerNotification(event, user) {
  await transporter.sendMail({
    from: '"Andon SEWS" <andon@sews.ma>',
    to: user.email,
    subject: `Alerte panne - Ligne ${event.ligne_id}`,
    html: `
      <h2>Alerte Andon</h2>
      <p>Ligne: ${event.ligne_id}</p>
      <p>Type: ${event.type}</p>
      <p>Heure: ${event.date_debut}</p>
    `
  });
}
\end{lstlisting}

\section{Synth\`ese du chapitre}

Le serveur backend du syst\`eme Andon est impl\'ement\'e en Node.js avec Express. Il int\`egre la connexion MQTT, la persistance PostgreSQL, l'authentification JWT, les routes API REST et la diffusion temps r\'eel via Socket.IO. Le code est organis\'e de mani\`ere modulaire pour faciliter la maintenance et les \'evolutions futures.
